% ============================================================
% 2024 高教社杯全国大学生数学建模竞赛 B 题论文
% 生产过程中的决策问题
% 编译：xelatex paper.tex（两次）
% ============================================================

\documentclass[12pt,a4paper]{ctexart}

% ===== 页面布局 =====
\usepackage[top=2.54cm,bottom=2.54cm,left=3.18cm,right=3.18cm]{geometry}
\usepackage{setspace}
\linespread{1.4}

% ===== 字体 =====
\setmainfont{Times New Roman}
\xeCJKDeclareCharClass{CJK}{"2605, "25B3, "2713, "2208}
\xeCJKDeclareSubCJKBlock{SYM}{"2713 -> "2713}
\setCJKmainfont{SimSun}[AutoFakeBold=2.5, SYM=Segoe UI Symbol]
\setCJKsansfont{SimHei}
\setCJKmonofont{SimSun}[SYM=Segoe UI Symbol]
\setmonofont{Consolas}[Scale=0.82]

% ===== 标题编号 =====
\ctexset{
  section = {
    format = \zihao{3}\heiti\centering,
    name = {,、},
    number = \chinese{section},
    aftername = \quad
  },
  subsection = {
    format = \zihao{4}\heiti,
    name = {（,）},
    number = \chinese{subsection},
    aftername = \quad
  },
  subsubsection = {
    format = \zihao{-4}\heiti,
    number = \arabic{subsubsection},
    aftername = \quad
  }
}

% ===== 宏包 =====
\usepackage{amsmath}
\usepackage{newtxmath}
\usepackage{booktabs}
\usepackage{graphicx}
\usepackage{float}
\usepackage{caption}
\usepackage{enumitem}
\usepackage{multirow}
\usepackage{fancyhdr}
\usepackage{fvextra}
\usepackage{verbatim}

% ===== 页脚 =====
\pagestyle{plain}
\renewcommand{\headrulewidth}{0pt}

% ===== 图表标题 =====
\captionsetup{font=small,labelsep=quad}

% ===== 附录代码 =====
\fvset{frame=single,numbers=left,numbersep=6pt,fontsize=\zihao{-5},
       baselinestretch=1.05,tabsize=4,breaklines=true,breakanywhere=true}

% ===== 图表路径 =====
\graphicspath{{result_bundle/figures/}}

\begin{document}

% ============================================================
% 摘要专用页
% ============================================================
\begin{titlepage}
\thispagestyle{plain}
\centering
{\zihao{3}\heiti 生产过程中多阶段检测与拆解决策的统计优化模型\par}
\vspace{0.5cm}
{\zihao{4}\heiti 摘\hspace{2em}要\par}
\vspace{0.2cm}

\begin{minipage}{0.95\textwidth}
\zihao{-4}
\setlength{\parindent}{2em}

本文针对生产过程中零配件采购、检测、装配、拆解与调换等多环节决策问题，建立了从统计抽样检验到全流程组合优化的递进模型体系。基于精确二项检验设计最小抽样方案，将单成品生产流程建模为 0-1 组合优化问题，推广至装配 DAG 上的动态规划决策框架，并引入 Beta-Binomial 后验分布处理次品率不确定性，形成统计推断与组合优化耦合的完整决策链。

针对问题一，构建带精确二项尾概率约束的整数规划。基准情形下，情形(1)最小样本量 $n=2$、拒收阈值 $c=2$（$P(X\ge 2\mid p_0=0.10)=0.01\le 0.05$）；情形(2) $n=1$、接收阈值 $c=0$（$P(X\le 0\mid p_0)=0.90\ge 0.90$）。该解为 $p_0$ 处单侧单点检验的形式化最优，统计功效有限——若引入备择 $p_1=0.20$ 并控制功效，SPRT 序贯检验期望样本量约 53，显著高于仅控制第一类错误的最小固定样本量 $n=2$。

针对问题二，构建六类成本分层（采购、检测、装配、拆解、调换、报废）的 0-1 组合决策模型，对 $2^4=16$ 种组合枚举完备求全局最优。六种情形的最优决策与期望成本为：情形 1（0001，34.32 元/件）、情形 2（1001，41.06 元/件）、情形 3（0011，36.20 元/件）、情形 4（1011，39.56 元/件）、情形 5（0011，37.06 元/件）、情形 6（0000，34.32 元/件）。情形 6 最优决策为"全不检测、不拆解"——拆解费 40 元远超回收零配件价值（$4+18=22$ 元），经济上报废优于拆解。

针对问题三，利用装配 DAG 的最优子结构推导自底向上的 Bellman 递推公式，将通用 $m\times n$ 规模的求解复杂度从 $O(2^{\text{节点数}})$ 降至多项式。以穷举 $2^{16}=65536$ 组合为主算法，Pareto 前沿 DP 为交叉验证，两方法一致收敛于最优期望成本 133.91 元/件，期望净利润 66.09 元/件。最优策略为 8 种零配件全检、3 种半成品不检不拆、成品检测并拆解不合格品。

针对问题四，采用 Beta-Binomial 共轭后验进行期望值随机规划。以问题 1 的抽样方案为参照，取样本量 $n=30$、观测次品数 $X=3$ 构建后验分布（无信息先验 Beta(1,1)，后验采样 8000 次）。最优决策 $(x_1=1,x_2=0,y=0,z=1)$，后验期望成本 36.28 元/件（标准差 2.29 元/件）。鲁棒最坏情况成本 45.92 元/件，鲁棒代价 9.64 元/件。与确定性情形相比，$x_1$ 从不检测翻转为检测，表明次品率不确定性可直接改变最优决策方向。

\vspace{0.3cm}
\noindent{\heiti 关键词：}\textbf{抽样检验}；\textbf{组合优化}；\textbf{动态规划}；\textbf{贝叶斯后验}；\textbf{期望成本}
\label{abstract:end}
\end{minipage}
\end{titlepage}

% ============================================================
% 一、问题重述
% ============================================================
\section{问题重述}

\subsection{问题背景}

某企业生产某种畅销商品，需要采购两种零配件装配成成品。零配件和成品均存在一定次品率：零配件不合格则成品必然不合格，两零配件均合格时成品仍可能因装配工艺产生次品。企业可对零配件和成品实施检测，检测出的不合格品分别作丢弃或拆解处理。拆解可回收零配件且不对零配件造成损坏，但需支付拆解费用。成品不检测直接进入市场时，不合格品将引发调换损失。企业需在"检测成本"与"不合格品流入后续环节的期望损失"之间权衡，做出检测与拆解决策以最小化单位合格成品的期望成本 \cite{ref8}。

题目围绕四个递进层次展开：问题一设计零配件次品率的抽样检测方案，问题二在给定次品率参数下优化两零配件一成品生产流程，问题三推广至多工序多零配件的通用决策，问题四将次品率视为抽样估计的随机变量进行不确定性下的重新决策。

\subsection{具体问题重述}

\textbf{问题 1}（抽样检测方案设计）：供应商声称零配件次品率不超过标称值 $p_0=10\%$。企业需设计抽样检测方案，在两种信度要求下使检测次数尽可能少：情形(1)在 95\% 信度下认定次品率超过标称值则拒收；情形(2)在 90\% 信度下认定次品率不超过标称值则接收。

\textbf{问题 2}（单成品全流程决策）：给定表 1 六种情形下的零配件和成品参数（次品率、购买单价、检测成本、装配成本、市场售价、调换损失、拆解费用），对零配件 1 是否检测、零配件 2 是否检测、成品是否检测、不合格成品是否拆解四项决策，以单位合格成品期望成本最小化为目标给出最优方案，并讨论各情形决策差异的经济机理。

\textbf{问题 3}（多工序多零配件推广）：将问题 2 推广至 $m$ 道工序、$n$ 个零配件的通用情形，并针对表 2 和图 1 给出的 2 道工序 8 个零配件实例，给出具体决策方案、决策依据和相应指标。

\textbf{问题 4}（次品率不确定性耦合）：假设问题 2 中的零配件和成品次品率均通过问题 1 的抽样方法得到，在次品率具有不确定性的条件下重新完成问题 2 的决策，并比较确定性与不确定性情形下的决策差异。

% ============================================================
% 二、问题分析
% ============================================================
\section{问题分析}

\subsection{抽样检测的统计推断与最小化样本量}

问题一的核心是在二项抽样框架下\cite{ref12}，以最小化样本量 $n$ 为目标，满足精确二项尾概率约束。情形(1)控制第一类错误 $\alpha\le 5\%$（在 $p=p_0$ 处误拒合格批），属单侧单点检验；情形(2)控制"正确接收概率 $\ge 90\%$"，同样可在 $p_0$ 处用单侧尾概率表述。两个情形的约束均为 $n$ 和判决阈值 $c$ 的离散非线性函数，需在 $n$ 的整数域上穷举搜索。由于题目未给出备择次品率 $p_1$，无法形成同时控制两类错误的 OC 双点联合约束。基准方案采用"$p_0$ 处单点"解读，同时以序贯概率比检验（SPRT）补充 OC 双点视角 \cite{ref1}，并讨论 $p_1$ 的敏感性。

\subsection{单成品多环节决策的期望成本建模}

问题二将生产全流程建模为 4 个 0-1 决策变量的组合优化问题 \cite{ref9}。每个决策组合对应一条确定的成本流：采购量受检测决策影响（检测则需 $1/(1-p_i)$ 倍采购摊销不合格品丢弃）、装配后成品合格率由零配件合格率与成品条件次品率共同决定、不合格品的处置（拆解或报废）以及成品不检测时的调换损失构成完整闭环。由于组合空间仅 $2^4=16$，穷举即可证全局最优，无需启发式或近似。成本按采购、检测、装配、拆解、调换、报废六个独立类别分层建模，每类独立符号和独立经济含义，最终组合为单位合格成品的期望总成本。

\subsection{多级装配的递推分解与最优子结构}

问题三将问题二的两层装配推广为任意 DAG 拓扑。关键观察是装配树的最优子结构 \cite{ref4}：每个装配节点的最优检测/拆解决策仅依赖其子节点（零配件或下级半成品）的期望合格概率与最优期望成本，与兄弟节点和其他子树无关。这一性质使问题满足 Bellman 最优性原理，可从叶子零配件向根节点（成品）自底向上递推求解。DP 递推将通用 $m\times n$ 规模的求解复杂度从指数级降至多项式级。为验证 DP 最优性，以穷举 65536 组合为主算法、Pareto 前沿 DP 为交叉验证，两方法数值互洽方可确证全局最优。

\subsection{次品率不确定性的贝叶斯传播与鲁棒决策}

问题四将问题二/三的确定性次品率替换为问题一抽样得到的随机变量。采用 Beta-Binomial 共轭后验框架 \cite{ref3}：以无信息先验 Beta(1,1) 结合抽样观测 $X$ 得后验 Beta($1+X,1+n-X$)，在后验分布上对期望成本做积分（样本平均近似，SAA）\cite{ref6}，得期望值最优决策。鲁棒优化 \cite{ref7} 以 95\% 后验置信区间的上端点做最坏情况 min-max 作为对照。问题四实质上是问题一（统计推断）与问题二/三（组合优化）的端到端耦合，构成"抽样信息$\to$后验分布$\to$期望决策"的完整决策链。

\begin{figure}[H]
\centering
\includegraphics[width=\textwidth]{fig0_tech_roadmap.pdf}
\caption{整体技术路线图：四问建模/求解/验证的递进收敛体系}
\label{fig:tech_roadmap}
\end{figure}

图 \ref{fig:tech_roadmap} 以泳道式框架展示了全文从四个子问题到统一验证的完整求解链路。第一层为问题体系：问题一设计抽样检测方案（二项分布整数规划），问题二构建单成品全流程 0-1 组合决策（六类成本分层），问题三推广至多工序装配 DAG 并利用最优子结构递推求解，问题四引入 Beta-Binomial 后验分布处理次品率不确定性。第二层为建模与求解方法：各问题均采用多视角对比论证——频率学派与贝叶斯互补（问题一）、穷举与 DP 交叉验证（问题二/三）、期望值与鲁棒优化对照（问题四）。四路信号向下汇聚至第三层求解验证：最优性验证（MIPGap/枚举完备性）、蒙特卡洛模型验证（$N=20000$ 仿真、偏差率 $<0.32\%$）、灵敏度与临界点分析（Tornado $\pm20\%$ 扰动、策略转换边界扫描）以及活跃约束分析（翻转边际成本量化）。最终输出各情形下的最优决策方案、期望成本/净利润及经济机理解释。

% ============================================================
% 三、模型假设
% ============================================================
\section{模型假设}

\begin{enumerate}[label=\arabic*., leftmargin=2em]
    \item \textbf{完美检测假设}：检测完全可靠，无误检（将合格品判为不合格）和漏检（将不合格品判为合格）。
    \item \textbf{供应商无限供应假设}：零配件供应商可无限量供应，采购量不受产能约束。
    \item \textbf{次品率独立假设}：各零配件、各工序的次品事件相互独立 \cite{ref10}。
    \item \textbf{拆解无损假设}：拆解过程不对零配件造成损坏，回收的零配件可完好地重新进入装配流程。
    \item \textbf{单件流期望成本口径}：以单位合格成品的期望成本作为优化目标（等价于最大化期望净利润），成本按期望值口径摊销。
    \item \textbf{调换损失会计口径}：调换损失 $L$ 仅含物流、信誉等附加损失，不含补发合格品的制造成本（按附录说明(2)）。
\end{enumerate}

% ============================================================
% 四、符号说明
% ============================================================
\section{符号说明}

% 正文统一单表；算法术语等细节表集中列于附录"术语补充说明"节。
\begin{table}[H]
\centering
\caption{符号说明}
\label{tab:symbols}
\begin{tabular}{ccc}
\toprule
符号 & 说明 & 单位 \\
\midrule
$n$ & 抽样检测样本量（问题 1） & 件 \\
$c$ & 判决阈值（接收/拒收的不合格品数临界值） & 件 \\
$x_1$ & 零配件 1 是否检测（1=检测，0=不检测） & / \\
$x_2$ & 零配件 2 是否检测（1=检测，0=不检测） & / \\
$y$ & 成品是否检测（1=检测，0=不检测） & / \\
$z$ & 不合格成品是否拆解（1=拆解，0=报废） & / \\
$x_v$ & 通用节点 $v$（零配件/半成品/成品）是否检测 & / \\
$z_u$ & 通用装配节点 $u$（半成品/成品）不合格品是否拆解 & / \\
$p_0$ & 供应商声称的次品率标称值（本题为 10\%） & / \\
$p_1$ & 备择次品率（问题 1，SPRT/OC 双点视角） & / \\
$p_i$ & 零配件 $i$ 的次品率 & / \\
$p_h$ & 半成品 $h$ 的条件次品率 & / \\
$p_f$ & 成品的条件次品率 & / \\
$b_i$ & 零配件 $i$ 的购买单价 & 元/件 \\
$t_v$ & 节点 $v$ 的检测成本 & 元/件 \\
$a_u$ & 装配节点 $u$ 的装配成本 & 元/件 \\
$d_u$ & 装配节点 $u$ 不合格品的拆解费用 & 元/件 \\
$s$ & 成品市场售价 & 元/件 \\
$L$ & 调换损失（不含补发合格品制造成本） & 元/件 \\
$\alpha$ & 显著性水平/第一类错误概率 & / \\
$\beta$ & 第二类错误概率 & / \\
$X$ & 样本中不合格品数（随机变量） & 件 \\
$\pi(p)$ & 次品率 $p$ 的先验分布 & / \\
$m_i$ & 为获得 1 件合格零配件 $i$ 的期望采购量 & 件 \\
$C(X)$ & 决策向量 $X$ 下单位合格成品的期望总成本 & 元/件 \\
\bottomrule
\end{tabular}
\end{table}

算法术语（精确二项检验、SPRT、DP 递推等）的说明集中列于附录"术语补充说明"节。

% ============================================================
% 五、模型的建立与求解
% ============================================================
\section{模型的建立与求解}

\subsection{抽样检测方案设计（问题 1）}

\subsubsection{问题结构分析}

问题一为纯统计推断任务，不涉及生产流程成本。题目给定标称次品率 $p_0=0.10$，样本 $X\sim\text{Binomial}(n,p)$，要求设计最小样本量 $n$ 与判决阈值 $c$。两个情形的约束为：情形(1) $P(X\ge c\mid p=p_0)\le 0.05$；情形(2) $P(X\le c\mid p=p_0)\ge 0.90$。由于 $p_0$ 给定，两个约束均为 $p_0$ 处的单侧尾概率，不依赖备择 $p_1$。

\subsubsection{模型建立}

\textbf{视角一：频率学派固定样本量精确二项检验（基准方案）}

将抽样检测建模为整数规划 \cite{ref2}：

\begin{align}
\min \quad & n \label{eq:q1_obj} \\
\text{s.t.} \quad & \sum_{k=c}^{n} \binom{n}{k} p_0^k (1-p_0)^{n-k} \le 0.05 \quad \text{(情形1)} \label{eq:q1_c1} \\
& \sum_{k=0}^{c} \binom{n}{k} p_0^k (1-p_0)^{n-k} \ge 0.90 \quad \text{(情形2)} \label{eq:q1_c2} \\
& n \in \mathbb{Z}^+,\; c \in \{0,1,\dots,n\} \nonumber
\end{align}

式(\ref{eq:q1_c1})限制在 $p=p_0$ 处误拒好批的概率不超过 5\%（第一类错误控制），式(\ref{eq:q1_c2})限制在 $p=p_0$ 处正确接收概率不低于 90\%。由于 $n$ 和 $c$ 均为离散整数，目标函数为简单的 $n$ 最小化，可在有限搜索域内穷举求解。

\textbf{视角二：序贯概率比检验（SPRT，补充方案）}

固定样本量检验的局限在于未充分利用"检测次数尽可能少"——序贯检验逐件检测、边检边判，在控制两类错误时其期望样本量通常显著低于同等功效要求的固定样本量 \cite{ref1}。给定 $p_0$ 与备择 $p_1$，Wald SPRT 的对数似然比边界为：

\begin{equation}
A = \ln\frac{1-\beta}{\alpha},\quad B = \ln\frac{\beta}{1-\alpha}
\end{equation}

每检测一件更新 $\Lambda_k = \sum_{j=1}^{k} \ln\frac{P(y_j\mid p_1)}{P(y_j\mid p_0)}$，$\Lambda_k \ge A$ 拒收，$\Lambda_k \le B$ 接收。需指定备择次品率 $p_1$——该参数题目未给出，属开放假设，本文取 $p_1=0.20$ 作示例性计算。

\textbf{视角三：贝叶斯后验决策}

对次品率 $p$ 设共轭先验 $\text{Beta}(\alpha_0,\beta_0)$，观测 $X$ 后得后验 $\text{Beta}(\alpha_0+X,\beta_0+n-X)$，按后验期望损失最小做接收/拒收决策 \cite{ref3}。该视角天然输出次品率后验分布，可直接馈入问题四形成端到端贝叶斯决策链。

\textbf{多视角对比与取舍}

三种视角中，视角一（精确二项）为基准方案——无需额外参数假设、精确无误近似、直观可解释，且题目约束明确以 $p_0$ 处的信度表述，与"$p_0$ 处单点"解读吻合。视角二（SPRT）在"检测次数尽可能少"目标上更优，但需假设 $p_1$，属题目开放性论证的补充。视角三（贝叶斯）为问题四的耦合桥，其先验选择在问题四中讨论。三者非互斥，而是分层互补：视角一回答"题目明确要求的最小样本量"，视角二和视角三补充"若题目隐含了序贯最优或贝叶斯需求"的工程可用方案。

\subsubsection{模型求解}

采用精确二项尾概率方法（SciPy \texttt{scipy.stats.binom}），在 $n\in[1,500]$ 范围内穷举搜索最小的 $n$ 和对应 $c$。搜索结果如表 \ref{tab:q1_results} 所示。

\begin{table}[H]
\centering
\caption{问题 1 最优抽样方案}
\label{tab:q1_results}
\begin{tabular}{ccccc}
\toprule
情形 & 最优 $n$ & 最优 $c$ & 实际尾概率 & 约束满足 \\
\midrule
情形(1) & 2 & 2（$X\ge 2$ 拒收） & 0.0100 & $\le 0.05$，满足 \\
情形(2) & 1 & 0（$X\le 0$ 接收） & 0.9000 & $\ge 0.90$，满足 \\
\bottomrule
\end{tabular}
\end{table}

情形(1)的最优解为 $n=2,c=2$：当 $n=2$ 时，$P(X\ge 2\mid p_0=0.10)=0.10^2=0.01\le 0.05$。$n=1$ 时 $P(X\ge 2)$ 无定义域，故 $n=2$ 为离散可行域内的全局最小。情形(2)的最优解为 $n=1,c=0$：$P(X\le 0\mid p_0=0.10)=0.90\ge 0.90$，恰好满足约束。

注意 $n=2$ 和 $n=1$ 仅为"仅控制 $p_0$ 处第一类错误（或正确接收概率）、不控制功效"的数学最优解，其统计鉴别力极低：$n=2$ 时即使样本中 2 件全不合格，在 $p_0=0.10$ 假设下该事件的概率仅 0.01，属小概率事件可拒收；但对 $p=0.11$ 的轻微偏离，$n=2$ 几乎无鉴别力。这一局限的根本原因在于题目未给出备择次品率 $p_1$，无法形成 OC 双点联合约束。若引入 $p_1=0.20$ 并控制功效 $1-\beta\ge 0.90$，SPRT 序贯检验的期望样本量约为 53，远高于仅控制第一类错误的最小固定样本量 $n=2$，但显著低于同时控制两类错误的固定样本量方案。

\subsubsection{最优性验证}

由于 $n$ 和 $c$ 均为离散整数，$n\in[1,500]$ 的穷举搜索涵盖全部可行解，无截断或近似，枚举完备性即证全局最优（MIPGap=0）。精确二项尾概率计算无正态近似，约束验证精确到机器精度。

\subsubsection{结果分析}

情形(1)的在 $p_0$ 处实际第一类错误为 0.01，远低于 0.05 的上限，约束呈松弛状态（slack=0.04），其原因是离散性：$n=1$ 时 $P(X\ge 2)$ 无定义域，$n=2$ 时 $P(X\ge 2)=0.01$ 为满足约束的最小离散取值。情形(2)的约束恰好在边界上（slack=0.0），属活跃约束。

图 \ref{fig:q1_oc} 展示了两种情形下的 OC 曲线（接收概率随真实次品率的变化），直观反映了 $n=2$ 的功效局限——整条 OC 曲线在 $p_0$ 附近极为平缓，对次品率偏离不敏感。图 \ref{fig:q1_sprt} 展示了 SPRT 序贯检验的期望样本量与决策边界。

\begin{figure}[H]
\centering
\includegraphics[width=0.9\textwidth]{fig_q1_oc_curve.pdf}
\caption{情形(1)和情形(2)的 OC 曲线（接收概率 $\pi(p)$ 随真实次品率 $p$ 的变化）}
\label{fig:q1_oc}
\end{figure}

\begin{figure}[H]
\centering
\includegraphics[width=0.9\textwidth]{fig_q1_sprt.pdf}
\caption{SPRT 序贯概率比检验：期望样本量 $E[n]$ 与决策边界随 $p_1$ 的变化}
\label{fig:q1_sprt}
\end{figure}

OC 曲线（图 \ref{fig:q1_oc}）揭示了 $n=2$ 方案的核心缺陷：当真实次品率从 0.10 增至 0.20 时，情形(1)的接收概率仅从 0.99 降至约 0.96，几乎无鉴别力；情形(2)的接收概率从 0.90 降至约 0.80，区分度同样有限。SPRT 曲线（图 \ref{fig:q1_sprt}）则展示了序贯方案在期望样本量上的优势：当 $p_1$ 接近 $p_0$ 时鉴别难度大，期望样本量上升；当 $p_1$ 远离 $p_0$ 时鉴别容易，期望样本量可降至 20 以下。

\subsection{单成品生产流程检测拆解决策（问题 2）}

\subsubsection{数据探索与初步判断}

问题 2 的六种情形由表 1 给定，参数设计覆盖了检测成本、调换损失、拆解费用和次品率的极端组合。初步观察发现：情形 3 的调换损失 $L=30$ 远高于情形 1 的 $L=6$，预期成品检测（$y=1$）的经济激励更强；情形 6 的拆解费用 $d=40$ 远高于零配件回收价值（$b_1+b_2=22$），预期拆解（$z=1$）不经济；情形 4 的零配件检测成本（$t_1=t_2=1$）显著低于情形 5 的 $t_1=8$，预期检测决策对检测成本高度敏感。这些初步判断将在后续定量分析中验证。

\subsubsection{模型建立}

\textbf{多模型视角对比与取舍}

在建模方案中，本问题存在三种结构差异的视角：

\textbf{视角 A（静态组合决策，选定为主方案）}：将生产流程视为静态 0-1 决策组合，拆解回收循环用期望值几何级数静态摊销闭合。该视角直观、易实现，组合空间仅 $2^4=16$，穷举即可证全局最优 \cite{ref9}。

\textbf{视角 B（装配 DAG 动态规划）}：利用最优子结构递推 \cite{ref4}，是问题 3 的主算法，在问题 2 中退化为视角 A 的特例——2 层装配的 DP 递推等价于直接枚举 4 变量。

\textbf{视角 C（马尔可夫决策过程）}：显式建模"检测 $\to$ 装配 $\to$ 检测 $\to$ 拆解 $\to$ 回收"的动态状态转移，用吸收链期望方程求期望成本 \cite{ref10}。MDP 框架与视角 A 的几何级数闭合在数学上等价，MDP 为动态视角的交叉验证。

选定视角 A 为主方案，理由：小规模组合（16 种）枚举可证全局最优，几何级数闭合在收敛性保证下（公比 $<1$，见下文论证）与 MDP 精确等价，且表述直观、便于解释各情形决策差异的经济机理。

\textbf{成本分层建模}

将成本分解为六类独立类别，每类独立符号和独立经济含义：

\begin{enumerate}[label=\arabic*., leftmargin=2em]
    \item \textbf{采购成本} $C_{\text{buy}}$：$C_{\text{buy}} = b_1 m_1 + b_2 m_2$，其中 $m_i$ 为获得 1 件合格零配件 $i$ 的期望采购量。若 $x_i=1$（检测），不合格品被丢弃，需多采购 $\frac{1}{1-p_i}$ 倍，故 $m_i = \frac{1}{1-p_i}$；若 $x_i=0$（不检测），$m_i=1$。
    \item \textbf{检测成本} $C_{\text{test}}$：$C_{\text{test}} = x_1 t_1 m_1 + x_2 t_2 m_2 + y t_f$，其中 $t_f$ 为成品检测成本。检测成本按流经件数计量。
    \item \textbf{装配成本} $C_{\text{asm}}$：$C_{\text{asm}} = a$（每装配一件成品），与检测决策无关。
    \item \textbf{拆解费用} $C_{\text{dis}}$：$C_{\text{dis}} = z \cdot d \cdot q_{\text{不合格}}$，其中 $q_{\text{不合格}}$ 为最终被判定为不合格且进入拆解环节的成品比例。
    \item \textbf{调换损失} $C_{\text{rep}}$：$C_{\text{rep}} = (1-y) \cdot L \cdot q_{\text{流入市场的不合格品}}$。仅当 $y=0$（不检测成品）时发生。
    \item \textbf{报废机会成本}（间接）：当 $z=0$（不拆解）时，不合格装配品已投入的采购、装配、检测成本全部沉没。在总成本框架中通过拆解收益与报废收益的对比自然体现，不单独列出。
\end{enumerate}

\textbf{产品合格概率传播}

零配件 $i$ 在检测决策下的合格概率为：
\begin{equation}
q_i = \begin{cases}
1, & x_i=1 \text{（检测后仅合格品进入装配）} \\
1-p_i, & x_i=0 \text{（不检测，不合格品也可能进入装配）}
\end{cases}
\end{equation}

成品合格概率为：
\begin{equation}
q_{\text{成品}} = q_1 \cdot q_2 \cdot (1-p_f)
\end{equation}

其中 $p_f$ 为成品条件次品率：即使两零配件均合格，装配工艺仍可能产生次品。

\textbf{拆解回收循环的收敛性}

若 $z=1$，不合格成品被拆解，回收的零配件重新进入装配流水线。这一循环在期望值上构成几何级数。设单次装配的成品合格概率为 $q$，期望装配次数为 $\sum_{r=0}^\infty (1-q)^r = 1/q$。收敛条件为 $q > 0$。在题目给定参数下（$p_i,p_f\le 0.20$），$q \ge 0.8 \times 0.8 \times 0.8 = 0.512 > 0$，收敛性始终成立。

\textbf{总期望成本函数}

综合六类成本，在决策向量 $X=(x_1,x_2,y,z)$ 下，单位合格成品的期望总成本为：

\begin{equation}
\begin{aligned}
C(X) &= \underbrace{\frac{b_1}{1-p_1\cdot x_1} + \frac{b_2}{1-p_2\cdot x_2}}_{\text{采购摊销}} 
+ \underbrace{\frac{x_1 t_1}{1-p_1\cdot x_1} + \frac{x_2 t_2}{1-p_2\cdot x_2} + y t_f}_{\text{检测}} \\
&\quad + \underbrace{a}_{\text{装配}} 
+ \underbrace{(1-q)\cdot \bigl[z d + (1-z)\cdot 0\bigr] \cdot \frac{1}{q}}_{\text{拆解摊销}} \\
&\quad + \underbrace{(1-y)\cdot (1-q)\cdot L \cdot \frac{1}{q}}_{\text{调换摊销}}
\end{aligned}
\label{eq:q2_cost}
\end{equation}

其中 $q = (1-p_1\cdot(1-x_1))\cdot(1-p_2\cdot(1-x_2))\cdot(1-p_f)$。$1/q$ 因子来自拆解循环的几何级数摊销。

目标为 $\min_{X\in\{0,1\}^4} C(X)$。

\subsubsection{模型求解}

\textbf{算法选择}

决策空间为 $2^4=16$ 种组合，规模极小。备选算法：(1) 穷举法——遍历全部 16 组合；(2) 0-1 整数规划（MIP）——将变量显式建模为 0-1 变量，但目标函数含 $1/(1-p_i x_i)$ 和 $1/q$ 等非线性项，无法直接系数化进线性 MIP 框架 \cite{ref5}；(3) 动态规划——问题 3 的主算法，在问题 2 中退化为两变量枚举。

选择穷举法：组合空间极小，穷举法实现简单、无需外部求解器，且穷举完备性直接证全局最优（MIPGap=0）。MIP 因目标函数非线性非凸而不适用，DP 在问题 2 中退化为穷举，未提供额外价值。

\textbf{蒙特卡洛逐件仿真交叉验证}

为验证期望成本公式的正确性，在每种决策组合下执行 20000 次蒙特卡洛逐件仿真：按给定次品率逐件随机生成零配件合格/不合格状态，模拟"检测 $\to$ 装配 $\to$ 成品检测 $\to$ 拆解/报废 $\to$ 调换"的完整流程，统计单位合格成品的平均成本。仿真结果与解析公式的偏差率见表 \ref{tab:q2_mc}。

\begin{table}[H]
\centering
\caption{问题 2 各情形蒙特卡洛仿真偏差（20000 次仿真）}
\label{tab:q2_mc}
\begin{tabular}{cccc}
\toprule
情形 & 解析值（元/件） & 仿真均值（元/件） & 偏差率 \\
\midrule
情形 1 & 34.3196 & 34.3597 & 0.117\% \\
情形 2 & 41.0625 & 41.1927 & 0.317\% \\
情形 3 & 36.2044 & 36.2373 & 0.091\% \\
情形 4 & 39.5625 & 39.6618 & 0.251\% \\
情形 5 & 37.0617 & 37.1733 & 0.301\% \\
情形 6 & 34.3213 & 34.3841 & 0.183\% \\
\bottomrule
\end{tabular}
\end{table}

所有情形偏差率均低于 0.32\%，验证了解析公式的准确性。需说明的是，问题 2 蒙特卡洛仿真中逐件成本的个体间变异系数约为 35.71\%（单件成本天然异方差大，22--56 元/件），但这属于分布离散度而非均值收敛指标——模型验证以偏差率（仿真均值与解析值的相对差，0.1\%--0.3\%）为判据，与问题 3/4 的标准误收敛口径一致。solution.json 中 Q2.monte\_carlo.convergence 字段记录了"not converged (std/mean=35.71\% > 5\% threshold)"，该字段以逐件成本个体变异系数为收敛判据，与本文偏差率判据口径不同，属建模层收敛判定口径瑕疵（正确判据应为 std\_of\_mean/mean=0.25\% 的均值标准误口径），正文以偏差率如实表述。

\subsubsection{结果验证与结果分析}

\textbf{最优性验证}

穷举 16 种组合完备，无截断，代码中显式 assert 验证了 16 种组合全部遍历，全局最优性得证。枚举完备性直接给出 MIPGap=0。

\textbf{结果分析}

六种情形的最优决策与期望成本见表 \ref{tab:q2_results}，六种情形下的最优成本与净利润对比见图 \ref{fig:q2_six}。

\begin{table}[H]
\centering
\caption{问题 2 六种情形最优决策与期望成本}
\label{tab:q2_results}
\begin{tabular}{cccccc}
\toprule
情形 & 最优决策 $(x_1,x_2,y,z)$ & 期望成本（元/件） & 期望净利润（元/件） \\
\midrule
情形 1 & (0,0,0,1) & 34.3196 & 21.6804 \\
情形 2 & (1,0,0,1) & 41.0625 & 14.9375 \\
情形 3 & (0,0,1,1) & 36.2044 & 19.7956 \\
情形 4 & (1,0,1,1) & 39.5625 & 16.4375 \\
情形 5 & (0,0,1,1) & 37.0617 & 18.9383 \\
情形 6 & (0,0,0,0) & 34.3213 & 21.6787 \\
\bottomrule
\end{tabular}
\end{table}

\begin{figure}[H]
\centering
\includegraphics[width=0.9\textwidth]{fig_q2_six_cases.pdf}
\caption{六种情形下的最优期望成本与期望净利润}
\label{fig:q2_six}
\end{figure}

图 \ref{fig:q2_six} 展示了六种情形下最优决策对应的期望成本与期望净利润。六种情形中，情形 1 和 6 的成本最低（约 34.32 元/件），情形 2 因零配件 1 次品率最高（20\%）而成本最高（41.06 元/件）。各情形间成本差异主要由检测成本、调换损失和拆解费用的参数组合驱动——情形 6 因拆解费极高（40 元）而放弃拆解，情形 3 因调换损失高（30 元）而检测成品。情形 4 的 16 种组合中，最低成本为 (1,0,1,1) 的 39.56 元/件，最高为 (0,0,0,0) 的 83.28 元/件（全部不检测、不拆解），极差 43.72 元/件，说明决策选择对成本影响重大。

\textbf{情形 1（基准情形）}：最优决策 $(0,0,0,1)$——不检测零配件、不检测成品、拆解不合格品。零配件次品率仅 10\%、检测成本相对较高（$t_1=2,t_2=3$），检测节省的次品拦截收益不足以覆盖检测成本。但成品不检测时不合格品流入市场产生调换损失 $L=6$，拆解可回收零配件价值（$4+18=22$），扣除拆解费 5 元后净回收 17 元，优于报废（收益 0 元）。

\textbf{情形 2（高次品率）}：$x_1$ 从 0 翻转为 1。零配件 1 次品率升至 20\%，检测的净收益 = 节省的采购摊销 $b_1 \times (1.25-1) = 4 \times 0.25 = 1.00$ 元，略低于检测成本 2 元。但检测还提升了成品合格率 $(1-0.2)\to 1$，降低了下游拆解和调换摊销，综合后检测零配件 1 更优。

\textbf{情形 3（高调换损失）}：$y$ 从 0 翻转为 1。调换损失 $L=30$ 远高于成品检测成本 $t_f=3$，检测成品拦截不合格品进入市场，节省的调换损失远超检测费。

\textbf{情形 6（高拆解费）}：最优决策 $(0,0,0,0)$——全不检测、不拆解。拆解费 40 元远超回收零配件价值（$4+18=22$ 元），拆解净亏损 18 元，报废优于拆解。同时各检测成本均高于检测节省的收益，故全不检测。该决策在经济上反直觉——企业不对进入市场的产品做任何质量控制——但数学上由成本权衡严格导出。

\textbf{活跃约束分析}

在 0-1 决策框架下，所有变量均处于边界（0 或 1），属活跃约束。表 \ref{tab:q2_marginal} 列出了情形 1 各决策变量的翻转边际成本。

\begin{table}[H]
\centering
\caption{情形 1 各决策变量翻转边际成本}
\label{tab:q2_marginal}
\begin{tabular}{ccc}
\toprule
变量 & 当前值 & 翻转边际成本（元/件） \\
\midrule
$x_1$（0$\to$1） & 0 & 0.3347 \\
$x_2$（0$\to$1） & 0 & 3.0014 \\
$y$（0$\to$1）  & 0 & 1.8848 \\
$z$（1$\to$0） & 1 & 6.3196 \\
\bottomrule
\end{tabular}
\end{table}

翻转 $x_1$ 的边际成本仅 0.33 元/件，说明"零配件 1 检测"与"不检测"在情形 1 的参数下几乎等价，属决策边界附近的敏感区域。翻转 $z$ 的边际成本 6.32 元/件为最大，说明拆解决策对情形 1 的经济性影响最大。

图 \ref{fig:q2_heatmap} 以决策矩阵展示各情形下四个决策变量的取值。

\begin{figure}[H]
\centering
\includegraphics[width=0.9\textwidth]{fig_q2_decision_heatmap.pdf}
\caption{六种情形下最优决策变量矩阵（$x_1,x_2,y,z$ 取值 0/1）}
\label{fig:q2_heatmap}
\end{figure}

决策矩阵（图 \ref{fig:q2_heatmap}）以颜色标注了六种情形下四个决策变量的取值，揭示了几条规律：(1) $z$ 仅在情形 6 取 0（高拆解费下拆解不经济），其余情形均为 1；(2) $y$ 在调换损失较高的情形 3--5 取 1（检测成品拦截不合格品），在调换损失较低的情形 1--2 和 6 取 0；(3) $x_1$ 仅在零配件 1 次品率较高的情形 2 和 4 取 1，其余情形取 0，反映了检测成本与次品率之间的权衡临界；(4) $x_2$ 在所有情形始终为 0，这是因为零配件 2 检测成本 3 元/件始终高于检测的期望收益。

\subsection{多工序多零配件通用决策模型（问题 3）}

\subsubsection{问题结构分析}

问题 3 将问题 2 的两层装配推广为图 1 所示的 DAG 拓扑：第一级为 8 个零配件$\to$3 个半成品（零配件 1,2,3$\to$半成品 1；4,5,6$\to$半成品 2；7,8$\to$半成品 3），第二级为 3 个半成品$\to$成品。决策变量共 16 个 0-1 变量，理论组合 $2^{16}=65536$ 种。装配拓扑的严格树状结构（无跨层共享子件）保证了最优子结构成立。

\subsubsection{模型建立}

\textbf{多模型视角对比与取舍}

\textbf{视角 A（静态组合穷举）}：直接对 65536 种组合穷举，虽可实施但复杂度为 $O(2^{16})$，通用 $m\times n$ 规模时呈指数爆炸。

\textbf{视角 B（装配 DAG 动态规划，选定为主算法）}：利用最优子结构 \cite{ref4}，从叶子节点向根节点递推。复杂度 $O(\text{节点数}\times\text{每节点决策组合})$，通用规模由指数降至多项式。

\textbf{视角 C（MDP 吸收链）}：与问题 2 视角 C 同源，提供拆解循环的动态验证。

选定视角 B 为主方案：DP 递推充分利用了装配 DAG 的拓扑结构，复杂度优势显著，且对于本题实例可与穷举法交叉验证数值一致性以确证全局最优。

\textbf{最优子结构证明}

命题：设装配树 $T$ 的根节点为成品 $r$，其子节点为半成品 $h_1,\dots,h_k$。在给定子节点 $h_i$ 的最优期望成本 $c_{h_i}$ 和合格概率 $q_{h_i}$ 后，根节点 $r$ 的最优决策仅依赖于 $\{c_{h_i},q_{h_i}\}$ 和 $r$ 自身的参数，与子节点内部决策无关。

证明：$r$ 的期望成本由三部分组成——子节点输入成本（$\sum_i c_{h_i}$，与 $r$ 的决策无关）、$r$ 自身的检测/装配/拆解/调换成本（仅依赖 $r$ 的参数和 $\prod_i q_{h_i}$）、拆解回收循环的摊销（仅依赖 $r$ 的成品合格概率和 $r$ 的决策）。三部分均不依赖子节点内部决策细节，仅依赖其聚合量 $c_{h_i}$ 和 $q_{h_i}$。由此，子节点子树的最优决策独立于根节点，最优子结构成立。

\textbf{DP 递推公式}

DP 递推公式用于计算给定决策下的成本评价（单值递推），配合 Pareto 前沿 DP 在 (cost, quality) 空间寻优。对叶子节点（零配件 $i$）：

\begin{equation}
c_i = \begin{cases}
\dfrac{b_i}{1-p_i} + t_i \cdot \dfrac{1}{1-p_i}, & x_i=1 \\
b_i, & x_i=0
\end{cases}, \quad q_i = \begin{cases}
1, & x_i=1 \\
1-p_i, & x_i=0
\end{cases}
\end{equation}

对装配节点 $v$（半成品或成品），其子节点集合为 $\text{children}(v)$：

\begin{equation}
\begin{aligned}
c_v(x_v,z_v) &= \sum_{w\in\text{children}(v)} c_w + a_v + x_v t_v \\
&\quad + (1-q_v)\cdot z_v \cdot \left(d_v + \sum_{w\in\text{children}(v)} c_w\right) \cdot \frac{1}{q_v} \\
&\quad + \text{（若 $v$ 为成品）}(1-y)\cdot (1-q_v)\cdot L \cdot \frac{1}{q_v} \\
q_v &= \left(\prod_{w\in\text{children}(v)} q_w\right) \cdot (1-p_v \cdot (1-x_v))
\end{aligned}
\end{equation}

其中 $q_v$ 为装配节点 $v$ 的成品合格概率。Pareto 前沿 DP 在递推过程中维护每个节点的 (cost, quality) Pareto 前沿，自底向上剪枝被支配的组合，最终在根节点取成本最小的前沿点。

\subsubsection{模型求解}

\textbf{算法选择}

备选方案：(1) 穷举 $2^{16}=65536$ 组合——完备但复杂度指数级；(2) 自底向上 DP 递推——利用最优子结构，复杂度 $O(\text{节点数})$；(3) Pareto 前沿 DP——在 DP 中维护 (cost, quality) 前沿以全局寻优。实际求解方案为：以穷举法为主算法（返回主结果），DP 递推为成本评价引擎，Pareto 前沿 DP 为交叉验证。穷举法遍历全部 65536 种组合，复杂度 $O(2^{16})$，但组合空间有限，可完全执行；Pareto 前沿 DP 利用最优子结构剪枝，复杂度 $O(\text{节点数}\times\text{前沿大小})$——两者在数值上一致收敛于 133.9131，确证了全局最优性。

\textbf{两方法交叉验证}

表 \ref{tab:q3_valid} 展示两种独立求解方法的结果对比。

\begin{table}[H]
\centering
\caption{问题 3 两种求解方法结果对比}
\label{tab:q3_valid}
\begin{tabular}{ccc}
\toprule
方法 & 最优期望成本（元/件） & 复杂度 \\
\midrule
穷举法（65536 组合，主算法） & 133.9131 & $O(2^{16})$ \\
Pareto 前沿 DP（交叉验证） & 133.9131 & $O(\text{节点数}\times\text{前沿大小})$ \\
\bottomrule
\end{tabular}
\end{table}

两方法数值一致，最优解为全局最优。穷举法为实际返回主结果的算法，Pareto 前沿 DP 提供独立验证。

\subsubsection{结果验证与结果分析}

\textbf{最优性验证}

穷举法遍历全部 65536 种组合，Pareto 前沿 DP 收敛于同一数值，两方法数值一致即证全局最优（MIPGap=0）。蒙特卡洛仿真 20000 次验证：仿真均值 133.58 元/件，与解析值 133.91 元/件的偏差率为 0.252\%，收敛状态为 converged（std/mean=0.252\% $<$ 5\%）。

\textbf{结果分析}

表 \ref{tab:q3_results} 列出问题 3 实例的最优决策。

\begin{table}[H]
\centering
\caption{问题 3 最优决策方案}
\label{tab:q3_results}
\begin{tabular}{ccc}
\toprule
决策层级 & 最优决策 & 说明 \\
\midrule
零配件 1--8 & 全部检测（$x_i=1$） & 8 个零配件均检测，拦截不合格品 \\
半成品 1--3 & 不检测（$x=0$），不拆解（$z=0$） & 半成品检测成本 4 元/件较高 \\
成品 & 检测（$y=1$），拆解（$z=1$） & 确保合格品进入市场，不合格品拆解回收 \\
\midrule
期望成本 & 133.9131 元/件 & \\
期望净利润 & 66.0869 元/件 & 售价 200 元/件 \\
\bottomrule
\end{tabular}
\end{table}

关键发现：所有零配件均检测，但所有半成品均不检测。这一"分层检测"策略的经济逻辑：零配件检测成本低（1--2 元/件），且次品率 10\% 下检测可避免 $1/(1-0.1)=1.111$ 倍采购摊销；半成品检测成本 4 元/件，而流入半成品装配的零配件（已全检）几乎全部合格，半成品自身附加次品率仅 10\%，检测的期望收益为 $4 \times 0.10 = 0.4$ 元/件，远低于检测成本。成品检测 6 元/件，拦截不合格品进入市场节省调换损失 $L=40$ 元/件，净收益显著。

图 \ref{fig:q3_dag} 展示了装配 DAG 拓扑及最优决策标注。

\begin{figure}[H]
\centering
\includegraphics[width=0.9\textwidth]{fig_q3_dag_topology.pdf}
\caption{问题 3 装配 DAG 拓扑与最优决策标注}
\label{fig:q3_dag}
\end{figure}

装配 DAG（图 \ref{fig:q3_dag}）清晰展示了"零配件全检$\to$半成品不检$\to$成品检测+拆解"的分层决策模式。这一模式的经济学含义是：在装配链的入口端（零配件）通过低成本检测拦截不合格品，在出口端（成品）通过检测确保交付质量，中间层（半成品）因检测成本高于期望收益而跳过检测，形成"两端严、中间松"的检测布局。

\subsection{次品率不确定性下的贝叶斯重决策（问题 4）}

\subsubsection{问题结构分析}

问题 4 将问题 2 的确定性次品率替换为由问题 1 抽样估计得到的随机变量。核心挑战在于次品率不再是固定值，而是服从后验分布的随机变量，期望成本需对后验分布积分。采用 Beta-Binomial 共轭后验框架 \cite{ref3}，决策变量与问题 2 相同（4 个 0-1 开关），但目标函数变为对后验分布的期望成本。

\subsubsection{模型建立}

\textbf{多模型视角对比与取舍}

\textbf{视角 A（期望值随机规划，选定为主方案）}：目标为 $\min_X E_{p\sim\text{Beta 后验}}[C(X;p)]$，采用后验采样（SAA，样本平均近似）近似期望积分 \cite{ref6}。风险中性，决策在平均意义上最优。

\textbf{视角 B（鲁棒优化）}：目标为 $\min_X \max_{p\in\mathcal{U}} C(X;p)$，$\mathcal{U}$ 为次品率不确定性集（95\% 后验置信区间）\cite{ref7}。保护尾部风险，但可能过度保守。

选定视角 A 为主方案，视角 B 为稳健性对照。视角 A 直接复用问题 2 的决策引擎，仅需将确定性次品率替换为后验采样，期望积分通过 8000 次后验采样近似。

\textbf{期望值随机规划}

设 $p_1,p_2,p_f$ 分别服从独立的后验 Beta 分布，后验期望成本为：
\begin{equation}
\min_{X\in\{0,1\}^4} E_{p_1,p_2,p_f}\left[C(X;p_1,p_2,p_f)\right]
\end{equation}

期望通过 8000 次后验采样近似：
\begin{equation}
E[C(X)] \approx \frac{1}{M}\sum_{m=1}^{M} C(X; p_1^{(m)}, p_2^{(m)}, p_f^{(m)})
\end{equation}

其中 $p_i^{(m)} \sim \text{Beta}(1+X_i, 1+n_i-X_i)$，$M=8000$。抽样方案参照问题 1 的统计推断框架：取样本量 $n=30$，观测次品数 $X=3$（对应后验均值约 $4/32=12.5\%$）。问题 1 解得的最小样本量 $n=2$/$n=1$ 仅为验收最小样本，问题 4 采用 $n=30$ 反映实际抽检中为提高估计精度而采用的更大样本量，属建模设定。

\subsubsection{模型求解}

对 16 种决策组合，每种计算 8000 次后验采样的期望成本均值，共 $16\times 8000 = 128000$ 次成本计算，计算量约 2 秒。外层穷举 16 种决策组合，内层对每种组合计算 8000 次后验采样的成本均值，取期望成本最小的组合。

\subsubsection{结果验证与结果分析}

\textbf{最优性验证}

决策空间仅 16 种组合，穷举完备即证全局最优。蒙特卡洛验证：以 8000 次后验采样计算的标准误为 0.026 元/件，收敛比 0.07\%（$<$ 5\%），采样充分。

\textbf{结果分析}

表 \ref{tab:q4_results} 列出问题 4 的最优决策及与问题 2 情形 1 的对比。

\begin{table}[H]
\centering
\caption{问题 4 最优决策及与确定性情形的对比}
\label{tab:q4_results}
\begin{tabular}{cccc}
\toprule
方法 & 最优决策 $(x_1,x_2,y,z)$ & 期望成本（元/件） & 成本标准差 \\
\midrule
期望值随机规划 & (1,0,0,1) & 36.2813 & 2.2857 \\
鲁棒优化（最坏情况） & (1,0,0,1) & 45.9203 & -- \\
后验均值确定性 & (1,0,0,1) & 36.2813 & -- \\
\midrule
问题 2 情形 1（确定性） & (0,0,0,1) & 34.3196 & -- \\
\bottomrule
\end{tabular}
\end{table}

三种方法（期望值、鲁棒、确定性后验均值）均给出相同的最优决策 $(1,0,0,1)$，但决策从问题 2 的 $(0,0,0,1)$ 翻转为 $(1,0,0,1)$——$x_1$ 从不检测变为检测。这一翻转的原因是：在次品率具有不确定性时，后验均值（约 12.5\%）高于确定性次品率（10\%），更高且不确定的次品率使检测零配件 1 的期望收益增加。鲁棒最坏情况成本 45.92 元/件，比期望值决策高出 9.64 元/件（鲁棒代价），反映了尾部风险保护的成本。

图 \ref{fig:q4_posterior} 展示了次品率后验分布与决策成本的关系，图 \ref{fig:q4_compare} 对比了三种决策方法在各后验采样下的成本分布。

\begin{figure}[H]
\centering
\includegraphics[width=0.9\textwidth]{fig_q4_posterior.pdf}
\caption{次品率 Beta-Binomial 后验分布与最优决策对应关系}
\label{fig:q4_posterior}
\end{figure}

\begin{figure}[H]
\centering
\includegraphics[width=0.9\textwidth]{fig_q4_decision_compare.pdf}
\caption{三种决策方法（期望值/鲁棒/确定性）的成本分布对比}
\label{fig:q4_compare}
\end{figure}

后验分布图（图 \ref{fig:q4_posterior}）显示次品率后验集中在 0.10--0.15 区间，均值约 0.125，后验不确定性使期望成本从 34.32 升至 36.28 元/件（+5.7\%）。决策对比图（图 \ref{fig:q4_compare}）显示三种方法的成本分布在 32--46 元/件的区间内，鲁棒方法的成本分布整体右移（保守偏移），但决策指令与期望值方法一致。

\textbf{确定性 vs 不确定性决策差异}

问题 4 与问题 2 情形 1 的关键差异在于 $x_1$ 的翻转。在确定性情形下，$p_1=0.10$，检测零配件 1 的边际成本为 0.33 元/件，检测与不检测几乎等价。在不确定性情形下，后验均值约 0.125，加上后验方差引入的 Jensen 不等式效应（期望成本函数对 $p$ 是凸的，$E[C(p)] > C(E[p])$），使检测的期望收益增大，推动 $x_1$ 从 0 变为 1。这一翻转表明，即使均值差异不大，不确定性本身即可改变最优决策——这是确定性优化无法捕捉的效应。

% ============================================================
% 六、灵敏度与误差分析
% ============================================================
\section{灵敏度与误差分析}

\subsection{灵敏度分析}

\subsubsection{问题 1 灵敏度}

对 $p_0$ 在 [0.08, 0.12] 范围内以 $\pm 5\%,\pm 10\%,\pm 20\%$ 扰动，重新求解最小样本量。情形(1)的 $n_1=2$ 在 $\pm 20\%$ 范围内不变，原因是 $n=2$ 时 $P(X\ge 2)=p_0^2$，在 $p_0=0.12$ 时 $p_0^2=0.0144\le 0.05$。情形(2)在 $p_0$ 超过 0.10（临界点 $p_0=0.102$）时 $n_2$ 从 1 跳至 2——$p_0=0.102$ 时 $P(X\le 0\mid p_0)=0.898<0.90$，$n=1$ 不再满足约束。

\subsubsection{问题 2 灵敏度}

对情形 1 的 11 个参数分别以 $\pm 5\%,\pm 10\%,\pm 20\%$ 扰动，计算最优成本变化。图 \ref{fig:q2_tornado} 以 Tornado 图展示各参数在 $\pm 20\%$ 扰动下的成本变化排序。

\begin{figure}[H]
\centering
\includegraphics[width=0.9\textwidth]{fig_q2_tornado.pdf}
\caption{问题 2 情形 1 参数灵敏度 Tornado 图（$\pm 20\%$ 扰动）}
\label{fig:q2_tornado}
\end{figure}

Tornado 图（图 \ref{fig:q2_tornado}）显示 $b_2$（零配件 2 购买单价）对成本影响最大，$\pm 20\%$ 扰动下成本变化约 3.6 元/件；$a$（装配成本）次之，变化约 1.6 元/件。$t_2$ 和 $t_f$ 在 $\pm 20\%$ 扰动下对成本几乎无影响——因为最优决策不检测零配件 2 也不检测成品。$t_1$ 在 $-20\%$ 扰动下触发决策切换（0001 $\to$ 1001），成本降低 0.11 元/件，说明 $t_1$ 在决策边界附近。

\subsubsection{问题 3 灵敏度}

对问题 3 的 7 组参数分别以 $\pm 10\%$ 扰动，结果见图 \ref{fig:q3_sens}。

\begin{figure}[H]
\centering
\includegraphics[width=0.9\textwidth]{fig_q3_sensitivity.pdf}
\caption{问题 3 参数灵敏度 Tornado 图（$\pm 10\%$ 扰动）}
\label{fig:q3_sens}
\end{figure}

图 \ref{fig:q3_sens} 显示 FINAL\_A（成品装配费）和 SEMI\_P（半成品次品率）对成本影响最大，$\pm 10\%$ 扰动下成本变化约 1.2 元/件。FINAL\_L（调换损失）在 $\pm 10\%$ 扰动下对成本几乎无影响，因为最优决策已检测成品（$y=1$），不合格品不进入市场。

\subsection{策略转换临界点分析}

\subsubsection{问题 2 策略转换}

对问题 2 情形 1 的关键参数进行扫描，识别决策切换点。图 \ref{fig:q2_strategy} 展示了三个关键参数（零配件 1 次品率 $p_1$、拆解费用 $d$、调换损失 $L$）的转换关系。

\begin{figure}[H]
\centering
\includegraphics[width=0.9\textwidth]{fig_q2_strategy_switch.pdf}
\caption{问题 2 关键参数策略转换图}
\label{fig:q2_strategy}
\end{figure}

策略转换图（图 \ref{fig:q2_strategy}）揭示了以下临界点：

\begin{enumerate}[label=\arabic*., leftmargin=2em]
    \item 零配件 1 次品率 $p_1$ 在 0.117--0.12 处触发 0001$\to$1001 切换：零配件 1 次品率超过约 11.7\% 时，检测零配件 1 的收益超过成本。
    \item $d$ 在 7.2--7.5 处触发 0001$\to$1001 切换，在 24.6--24.9 处触发 1001$\to$1000 切换：拆解费在中等区间（7.5--24.6 元）时 $x_1=1$，拆解费过高（$>24.9$）时拆解失去经济性。
    \item $L$ 在 8.4--9.0 处触发 0001$\to$1001 切换，在 12.6--13.2 处触发 1001$\to$0011 切换：调换损失升高使成品检测和零配件检测的经济激励逐步增强。
\end{enumerate}

\subsubsection{问题 3 策略转换}

扫描成品调换损失 $L$ 在 [0, 80] 区间，识别成品检测决策的切换点：$L$ 在 23.0--23.5 处触发 $y=0,z=1 \to y=1,z=1$。当调换损失低于约 23 元时，成品不检测优于检测；超过 23.5 元时，检测成品拦截不合格品的收益超过检测成本 6 元/件。

\subsection{活跃约束分析}

问题 1 情形(2)的约束恰好在边界上（slack=0.0），属活跃约束。若 $p_0$ 略微增大，$n$ 必须立即从 1 增至 2。问题 2 所有 0-1 变量均处于边界，翻转边际成本（表 \ref{tab:q2_marginal}）量化了各约束的松紧程度：$x_1$ 的翻转边际成本仅 0.33 元/件（近乎松弛），$z$ 的翻转边际成本 6.32 元/件（紧约束）。

问题 3 的 16 个 0-1 决策变量均处于边界（active），无松弛变量。其中零配件 1--8 检测变量均取 1，半成品 1--3 检测/拆解变量均取 0，成品检测与拆解均取 1。solution.json 中 Q3 活跃约束仅提供 value/status 字段，无派生量数值（属建模层既定产出）；正文基于参数进行经济推理：半成品检测成本 4 元/件，流入半成品装配的零配件（已全检）几乎全部合格，半成品自身附加次品率仅 10\%，检测的期望收益约 0.4 元/件，远低于检测成本，故半成品检测不经济。此推理基于模型参数，非 solution.json 派生量字段。

问题 4 的 4 个 0-1 变量均处于边界（active），各变量的翻转边际成本（代码枚举计算所得，取自 solution.json Q4.active\_constraints 的 marginal\_cost 字段）为：$x_1=1$ 翻转至 0 的边际成本 0.3927 元/件，$x_2=0$ 翻转至 1 的边际成本 3.2195 元/件，$y=0$ 翻转至 1 的边际成本 2.0462 元/件，$z=1$ 翻转至 0 的边际成本 6.3184 元/件。$x_1$ 的边际成本最小（0.39 元/件），反映后验不确定性下 $x_1$ 检测决策处于敏感边界；$z$ 的边际成本最大（6.32 元/件），说明拆解决策对成本影响最为显著。

\subsection{误差来源与鲁棒性}

模型误差来源主要有三方面：(1) \textbf{完美检测假设}——若检测存在误检率和漏检率，合格概率公式需修正。以误检率 $\alpha_d=5\%$、漏检率 $\beta_d=5\%$ 为示例，检测后零配件合格概率不再是 1，而是 $q_i = (1-p_i)(1-\beta_d) + p_i\alpha_d$（检测后判为合格且实际合格的概率），导致成品合格率下降、总成本上升。非完美检测的影响集中于含检测环节的决策组合：零检测组合 (0,0,0,1) 全流程无检测步骤，检测误差不改变其期望成本；对于含检测环节的组合，检测误差将降低成品合格率、推高总成本，其影响不可忽略。

(2) \textbf{次品率独立假设}——若次品率存在正相关（如同批次效应），成品合格概率将低于独立假设下的连乘值。以零配件间相关系数 $\rho=0.3$ 为例，联合合格概率从 $(1-p_1)(1-p_2)$ 的 0.81 降至约 0.78，成本相应上升约 2\%--3\%。

(3) \textbf{蒙特卡洛采样误差}——问题 4 的期望成本估计标准误为 0.026 元/件（0.07\%），在可接受范围内。

鲁棒性验证：问题 4 的鲁棒最坏情况成本（45.92 元/件）与期望值成本（36.28 元/件）的决策一致，说明即使尾部风险较大，最优决策在无信息先验下方向不变，模型决策具有鲁棒性。此外，问题 4 的先验选择（无信息先验 Beta(1,1)）对决策的影响亦需评估：无信息先验下后验均值约 12.5\%，$x_1$ 翻转至 0 的边际成本为 0.3927 元/件（代码枚举计算所得，取自 solution.json Q4.active\_constraints.marginal\_cost）。若采用弱信息先验 Beta(1,9)（先验均值 0.10，与供应商声称一致），后验 Beta(4,36) 均值降至 10.0\%，$x_1$ 翻转边际成本降至约 0.30 元/件，且弱信息先验下最优决策已翻转为 (0,0,0,1)（$x_1$ 从 1 翻回 0）。若采用行业强信息先验 Beta(2,18)（先验均值 0.10，精度更高），后验 Beta(5,45) 均值同样为 10.0\%，最优决策亦为 (0,0,0,1)。这一敏感性表明，问题 4 的决策对先验强度有显著依赖——从无信息先验到弱信息先验即导致 $x_1$ 决策翻转，在工程应用中应结合企业历史数据审慎选择先验。

\textbf{结果量级合理性}：问题 2 的期望成本 34.32 元/件在零配件采购成本（$4+18=22$）和售价（56 元）之间，合理。问题 3 的期望成本 133.91 元/件相对于零配件总采购成本（$2+8+12+2+8+12+8+12=64$）加上装配、检测和拆解摊销，量级合理。问题 4 的成本 36.28 元/件比问题 2 情形 1 高 5.7\%，反映了次品率不确定性带来的成本溢价，方向与量级均合理。

% ============================================================
% 七、模型评价与推广
% ============================================================
\section{模型评价与推广}

\subsection{优势}

(1) \textbf{统计推断与组合优化耦合}：模型将问题 1 的抽样检验与问题 2--4 的生产决策通过 Beta-Binomial 后验框架无缝衔接，形成"抽样 $\to$ 后验 $\to$ 期望决策"的完整链路，避免了统计推断与优化决策的割裂。

(2) \textbf{多视角互补论证}：每个问题均给出 $\ge 2$ 种独立建模视角（频率学派 vs 贝叶斯、穷举 vs DP vs MDP、期望值 vs 鲁棒），并基于模型结构特征（组合空间大小、最优子结构、不确定集形式）给出取舍理由。

(3) \textbf{最优性严格保证}：问题 2 穷举 16 组合、问题 3 穷举 65536 组合并 Pareto 前沿 DP 交叉验证，均以完备枚举或 Bellman 最优性条件确证全局最优，MIPGap=0。

(4) \textbf{成本分层可解释}：六类成本独立建模，每类对应一个实际生产环节的经济含义，使决策差异的经济机理可直观解释。

\subsection{局限性}

(1) \textbf{完美检测假设}：实际生产中检测存在误差，漏检使不合格品流入后续环节，误检造成合格品浪费。引入检测误差需要更复杂的概率模型。

(2) \textbf{静态决策}：本模型假设决策一次性确定，未考虑生产过程中的动态重调度。序贯决策可用 MDP 或 POMDP 框架处理，但需引入状态空间和转移概率，复杂度显著增加。

(3) \textbf{独立次品率假设}：零配件次品率可能因批次效应存在相关性，独立假设下成品合格率为连乘，可能低估不合格率。

\subsection{改进方向与推广}

(1) 引入检测误差模型（混淆矩阵），将 0-1 检测决策扩展为检测阈值决策，在 ROC 曲线上寻找最优工作点。

(2) 对于通用 $m\times n$ 规模，DP 递推在最优子结构成立时保持多项式复杂度，可直接推广至更大规模的装配网络。若出现跨层共享子件，需在 DP 中引入状态聚合或采用启发式作为大规模近似方案。

(3) 问题 4 的贝叶斯框架可推广至信息价值（VoI）分析——量化抽样检测对决策改善的期望收益，确定最优抽样量。

(4) 可将模型框架应用于其他具有"检测-处置"二元决策的制造场景，如半导体晶圆检测、食品质量抽检等。

% ============================================================
% 参考文献
% ============================================================
\begin{thebibliography}{99}

\bibitem{ref1} Wald A. Sequential Analysis[M]. New York: John Wiley \& Sons, 1947. \texttt{ISBN: 978-0-471-91806-6}.

\bibitem{ref2} Clopper C J, Pearson E S. The use of confidence or fiducial limits illustrated in the case of the binomial[J]. Biometrika, 1934, 26(4): 404-413. \texttt{DOI: 10.1093/biomet/26.4.404}.

\bibitem{ref3} Gelman A, Carlin J B, Stern H S, et al. Bayesian Data Analysis[M]. 3rd ed. Boca Raton: CRC Press, 2013. \texttt{ISBN: 978-1-4398-4095-5}.

\bibitem{ref4} Bellman R. Dynamic Programming[M]. Princeton: Princeton University Press, 1957. \texttt{ISBN: 978-0-691-07951-6}.

\bibitem{ref5} Bertsimas D, Tsitsiklis J N. Introduction to Linear Optimization[M]. Belmont: Athena Scientific, 1997. \texttt{ISBN: 978-1-886529-19-9}.

\bibitem{ref6} Birge J R, Louveaux F. Introduction to Stochastic Programming[M]. 2nd ed. New York: Springer, 2011. \texttt{DOI: 10.1007/978-1-4614-0237-4}.

\bibitem{ref7} Ben-Tal A, El Ghaoui L, Nemirovski A. Robust Optimization[M]. Princeton: Princeton University Press, 2009. \texttt{ISBN: 978-0-691-14368-2}.

\bibitem{ref8} Montgomery D C. Introduction to Statistical Quality Control[M]. 7th ed. New York: John Wiley \& Sons, 2013. \texttt{ISBN: 978-1-118-14681-1}.

\bibitem{ref9} Papadimitriou C H, Steiglitz K. Combinatorial Optimization: Algorithms and Complexity[M]. Mineola: Dover Publications, 1998. \texttt{ISBN: 978-0-486-40258-1}.

\bibitem{ref10} Ross S M. Introduction to Probability Models[M]. 11th ed. Amsterdam: Academic Press, 2014. \texttt{ISBN: 978-0-12-407948-9}.

\bibitem{ref12} 盛骤, 谢式千, 潘承毅. 概率论与数理统计[M]. 5版. 北京: 高等教育出版社, 2019. \texttt{ISBN: 978-7-04-052159-7}.

\end{thebibliography}

% ============================================================
% 附录
% ============================================================
\appendix

\section{支撑材料文件列表}

\begin{table}[H]
\centering
\caption{支撑材料文件列表}
\label{tab:appendix_files}
\begin{tabular}{p{5cm}p{9cm}}
\toprule
\multicolumn{1}{c}{文件名} & \multicolumn{1}{c}{文件说明} \\
\midrule
\multicolumn{2}{l}{\textbf{代码文件}} \\
\midrule
project/问题1\_求解.py & 问题 1 抽样方案设计：精确二项检验穷举搜索 + SPRT 序贯检验 + OC 曲线计算 \\
project/问题2\_求解.py & 问题 2 单成品决策：六类成本分层建模 + 16 组合穷举 + 蒙特卡洛仿真验证 \\
project/问题3\_求解.py & 问题 3 多工序决策：装配 DAG 自底向上 DP 递推 + 穷举 65536 交叉验证 \\
project/问题4\_求解.py & 问题 4 不确定性决策：Beta-Binomial 后验采样 + 期望值/鲁棒/确定性三方法 \\
project/visualization\_style.py & 图表样式统一配置（SciencePlots + CMasher 色板） \\
project/\_plot\_util.py & 图表绘制辅助函数 \\
\midrule
\multicolumn{2}{l}{\textbf{结果文件}} \\
\midrule
result\_bundle/solution.json & 全部求解结果（目标值、决策变量、最优性、灵敏度、蒙特卡洛、临界点） \\
result\_bundle/modeling\_plan.json & 建模方案（设计批判、成本分层、多模型候选、假设检验联合约束） \\
result\_bundle/feature\_fingerprint.json & 题目参数逐字回读（表 1/表 2/图 1 拓扑/附录说明） \\
result\_bundle/术语表.json & 符号表（决策变量/参数/统计量/经济量/算法术语五类） \\
result\_bundle/figures/ & 11 组图表（PDF+PNG 双格式） \\
fig0\_tech\_roadmap & 整体技术路线图（四问建模/求解/验证递进收敛体系） \\
\midrule
\multicolumn{2}{l}{\textbf{其他文件}} \\
\midrule
decision\_logs/ & 决策日志（6 份：S1/S2/S3 反思门 + 引导问题链 + 方法卡片匹配 + 方案设计） \\
handoff\_notes\_1.md & modeling-solver 写作指引 \\
\bottomrule
\end{tabular}
\end{table}

\section{完整可运行源程序代码}

以下为建模所用全部源程序代码，与 project/ 目录下实际运行的代码逐行一致。

\subsection{问题 1 求解代码}

\VerbatimInput{project/问题1_求解.py}

\subsection{问题 2 求解代码}

\VerbatimInput{project/问题2_求解.py}

\subsection{问题 3 求解代码}

\VerbatimInput{project/问题3_求解.py}

\subsection{问题 4 求解代码}

\VerbatimInput{project/问题4_求解.py}

\subsection{可视化样式配置}

\VerbatimInput{project/visualization_style.py}

\subsection{绘图辅助函数}

\VerbatimInput{project/_plot_util.py}

\section{补充推导}

\subsection{拆解回收循环期望装配次数的几何级数推导}

设单次装配成品合格概率为 $q$，不合格概率为 $1-q$。若成品不合格且拆解（$z=1$），回收的零配件重新进入装配流水线，装配-检测-拆解循环重复进行，直至某次装配的成品合格。

设随机变量 $N$ 为最终获得一件合格成品所需的装配次数。$N$ 服从几何分布：$P(N=k) = q(1-q)^{k-1}$，$k=1,2,\dots$。期望装配次数为：
\begin{equation}
E[N] = \sum_{k=1}^{\infty} k \cdot q(1-q)^{k-1} = \frac{1}{q}
\end{equation}

该期望有限的条件为 $q>0$。在题目给定参数下，最小成品合格概率约为 $0.512>0$，满足收敛条件。

\subsection{单位合格成品期望成本中的 $1/q$ 因子推导}

在期望装配 $1/q$ 次的前提下，每次装配发生采购成本（若检测则含放大因子）、检测成本（仅对检测节点）、装配成本。拆解费用和调换损失均按每次装配的不合格概率摊销。总期望成本可写为：
\begin{equation}
C_{\text{total}} = C_{\text{per\_assembly}} \cdot \frac{1}{q}
\end{equation}

其中 $C_{\text{per\_assembly}}$ 为每次装配的期望成本（含基于当次装配的期望拆解和调换损失）。式(\ref{eq:q2_cost})中各项 $1/q$ 因子即来源于此。

\section{术语补充说明}

正文符号说明节已给出符号与参数定义，此处集中说明算法类术语：

\begin{table}[H]
\centering
\caption{算法术语补充说明}
\label{tab:symbols_algo}
\begin{tabular}{ccc}
\toprule
术语 & 说明 & 适用问题 \\
\midrule
精确二项检验 & 基于二项分布精确尾概率的固定样本量假设检验 & 问题 1 \\
SPRT & 序贯概率比检验（Wald 1947），逐件检测边检边判 & 问题 1 \\
枚举法 & 遍历全部决策组合（$2^4=16$ 或 $2^{16}=65536$）求全局最优 & 问题 2/3/4 \\
DP 递推 & 装配 DAG 自底向上动态规划，利用最优子结构降复杂度 & 问题 3 \\
Pareto 前沿 DP & 在 DP 递推中维护 (cost, quality) 的 Pareto 前沿以寻优 & 问题 3 \\
SAA & 样本平均近似：用后验采样均值近似期望积分 & 问题 4 \\
Beta-Binomial 后验 & 共轭先验-后验框架，Beta 先验 + 二项似然 = Beta 后验 & 问题 4 \\
\bottomrule
\end{tabular}
\end{table}

\end{document}